% Source: https://artofproblemsolving.com/community/c3h1760131_is_there_a_combinatorical_reason_for_the_sum_of_the_first_n_cubes
% Stable identity: aops-topic-1760131-q01
Count the ordered pairs $(S,T)$ of $2$-element subsets of $\{1,2,\dots,n+1\}$. There are $\dbinom{n+1}{2}^2$ of them.

Now sort them by $d=\max(S\cup T)$, which can be $2,3,\dots,n+1$. Fix $d$; all elements of $S$ and $T$ lie in $\{1,\dots,d\}$ and $d$ is the largest element of at least one of them.
\begin{itemize}
\item $\max S=\max T=d$: each of $S,T$ is $d$ together with one of $1,\dots,d-1$: $(d-1)^2$ pairs.
\item $\max S=d>\max T$: $S$ in $d-1$ ways, $T\subseteq\{1,\dots,d-1\}$ in $\dbinom{d-1}{2}$ ways: $(d-1)\dbinom{d-1}2$ pairs; the case $\max T=d>\max S$ gives the same number.
\end{itemize}
So the number with maximum $d$ is
\[
(d-1)^2+2(d-1)\cdot\dfrac{(d-1)(d-2)}{2}=(d-1)^2\bigl(1+(d-2)\bigr)=(d-1)^3 .
\]
Summing over $d=2,\dots,n+1$ counts every pair once, so $\dbinom{n+1}{2}^2=\displaystyle\sum_{d=2}^{n+1}(d-1)^3=1^3+2^3+\cdots+n^3$. $\blacksquare$
